\section{Introducción}
Para la realización de este segundo proyecto de la materia de Modelado y Programación se nos encomendó la realización de un Reproductor de Música, el cual implementa una base de datos SQLite, misma que se nos fue proporcionada al inicio del mismo proyecto.
Las especificaciones de este programa, fueron varias, entre los que más destacan son:
\begin{enumerate}
    \item Que el reproductor sea capaz de leer las entradas de la base de datos que se nos fue proporcionada
    \item Los archivos de audio deben de tener la etiqueta de ID3V2.4, y con los metadatos que contenga el archivo se llenen los parámetros establecidos en la base de datos
    \item El archivo .db que se genere para la base de datos se almacene dentro del directorio estándar de datos de cada sistema operativo.
    \item A diferencia del primer proyecto, en donde era opcional la implementación de una interfaz gráfica, en este proyecto se exige la implementación de una.
    \item Dentro de la interfaz gráfica debe de ser capaz de poder editar las entradas de la base de datos, cabe aclarar que se especifica que los archivos de audio NO deben de ser modificados.
    \item La implementación de un ``minero'' el cual se encargará de escanear los archivos mp3, que sean ID3V2, y que además los inserte dentro de la base de datos, para que así, los pueda reproducir la interfaz gráfica
    \item Que se puedan reproducir los archivos dentro de la interfaz gráfica
    \item Cada uno de los componentes (Minero, Reproductor) deben de tener su propio ejecutable
\end{enumerate}
Con estos parámetros fue como se inicio con la parte del análisis y modelado del programa.